%!TEX program = xelatex
\def\numeromodulo{1}
\def\numerolista{3}
\documentclass[12pt]{exam}

\def\ano{2023}
\def\semestre{1}
\def\disciplina{Introdução à \'Algebra Linear}
\def\turma{3}

\usepackage{caption}
\usepackage{amssymb}
\usepackage{amsmath,amsfonts,amsthm,amstext}
\usepackage[utf8]{inputenc} % Required for inputting international characters
\usepackage[T1]{fontenc}
\usepackage[brazil]{babel}
\usepackage{graphicx}
\graphicspath{{../figuras_geogebra/}{../Pictures/}}
\usepackage{enumitem}
\usepackage{multicol}
\usepackage{answers}
\usepackage[svgnames]{xcolor}
\usepackage{tikz}
\usepackage{ifthen}
\usetikzlibrary{lindenmayersystems}
\usetikzlibrary[shadings]
\usepackage{tikz-3dplot}
\usepackage{tcolorbox}
%-------------amatrix
% Augmented matrix.  Usage (note the argument does not count the aug col):
% \begin{amatrix}{2}
%   1  2  3 \\  4  5  6
% \end{amatrix}
\newenvironment{amatrix}[1]{%
  \left(\begin{array}{@{}*{#1}{r}|r@{}}
}{%
  \end{array}\right)
}

\Newassociation{solucao}{Solution}{ans}
\newtheorem{exercicio}{}

\setlength{\topmargin}{-1.0in}
\setlength{\oddsidemargin}{0in}
\setlength{\textheight}{10.1in}
\setlength{\textwidth}{6.5in}
\setlength{\baselineskip}{12mm}

\newcounter{exercicios}
\setcounter{exercicios}{0}
\newcommand{\questao}{
\addtocounter{exercicios}{1}
\noindent{\bf Exerc{\'\i}cio \Roman{exercicios}: }}

\newcommand{\resp}[1]{
\noindent{\bf Exerc{\'\i}cio #1: }}

\extraheadheight{0.7in}
\firstpageheadrule
\runningheadrule
\lhead{
        \begin{minipage}[c]{1.7cm}
        \includegraphics[width=1.7cm]{unb.pdf}
        \end{minipage}%
        \hspace{0pt}
        \begin{minipage}[c]{4in}
          {Universidade de Bras{\'\i}lia} --
          {Departamento de Matem{\'a}tica}
\end{minipage}
\vspace*{-0.8cm}
}
% \chead{Universidade de Bras{\'\i}lia - Departamento de Matem{\'a}tica}
% \rhead{}
% \vspace*{-2cm}

\extrafootheight{.5in}
\footrule
\lfoot{\disciplina\ - \semestre$^o$/\ano\ - Módulo \numeromodulo}
\cfoot{}
\rfoot{P\'agina \thepage\ de \numpages}
\newcommand{\nsub}{\varsubsetneq}
\newcommand{\vaz}{\emptyset}
\newcommand{\im}{{\rm Im\,}}
\newcommand{\sub}{\subseteq}
\renewcommand{\sin}{{\rm sen\,}}
\newcommand{\n}{\mathbb{N}}
\newcommand{\z}{\mathbb{Z}}
\newcommand{\rac}{\mathbb{Q}}
\newcommand{\real}{\mathbb{R}}
\newcommand{\complex}{\mathbb{C}}
\newcommand{\cp}[1]{\mathbb{#1}}
\newcommand{\ch}{\mbox{\textrm{car\,}}\nobreak}
\newcommand{\vesp}[1]{\vspace{ #1  cm}}
\newcommand{\compcent}[1]{\vcenter{\hbox{$#1\circ$}}}
\newcommand{\comp}{\mathbin{\mathchoice
{\compcent\scriptstyle}{\compcent\scriptstyle}
{\compcent\scriptscriptstyle}{\compcent\scriptscriptstyle}}}

\newcommand{\integer}{\mathbb{Z}}
\begin{document}

  \Opensolutionfile{ans}[ans1]
  \begin{center}
    {\Large\bf \disciplina\ - Turma \turma\ -- \semestre$^{o}$/\ano} \\ \vspace{9pt} {\large\bf
        $\numerolista^a$ Lista de Exerc{\'\i}cios -- Módulo \numeromodulo}\\ \vspace{9pt} Prof. Jos{\'e} Ant{\^o}nio O. Freitas
  \end{center}
  \hrule

\vesp{.6}

\begin{exercicio}
  A figura abaixo mostra uma rede de canos de água com fluxo medido em litros por minuto.
  \begin{enumerate}
    \item Monte e resolva um sistema de equações lineares para encontrar os fluxos possíveis.
    \item Se o fluxo através de $AB$ é restrito a 5L/min, qual será o fluxo através dos outros dois ramos?
    \item Quais são os fluxos mínimo e máximos possíveis através de cada ramo?
  \end{enumerate}
  \begin{solucao}
    \begin{enumerate}
      \item $\begin{cases}f_1 = 30 - t\\f_2 = -10 + t\\f_3 = t\end{cases}$
      \item $f_1 = 15$, $f_3 = 15$
      \item $0 \le f_1 \le 20$, $0 \le f_2 \le 20$, $10 \le f_3 \le 30$
    \end{enumerate}
  \end{solucao}
\end{exercicio}

\begin{exercicio}
  A figura abaixo mostra uma rede viária de ruas de mão única com fluxo de tráfego nos sentidos indicados. As taxas
  de fluxo ao longo das ruas são medidas pelo número médio de veículos por hora.
  \begin{enumerate}
    \item Monte um sistema linear cuja solução forneça as taxas de fluxo desconhecidas.
    \item Resolva o sistema para as taxas de fluxo desconhecidas.
    \item Se o fluxo ao longo da rua de $A$ para $B$ precisar ser reduzido em virtude de uma obra, qual será o fluxo mínimo necessário
      para manter o tráfego fluindo em todas as ruas?
  \end{enumerate}
  \begin{solucao}
    \begin{enumerate}
      \item $\begin{cases}x_3 - x_4 = -500\\-x_1 + x_4 = 100\\x_1 - x_2 = 300\\x_2 - x_3 = 100\end{cases}$
      \item $x_1 = -100 +t$, $x_2 = -400 + t$, $x_3 = -500 + t$, $x_4 = t$
      \item $t = 500$
    \end{enumerate}
  \end{solucao}
\end{exercicio}

\end{document}
